\documentclass[10pt,a4paper]{article}

\usepackage[utf8]{inputenc}
\usepackage[T1]{fontenc}
\usepackage[margin=1.5cm]{geometry}
\usepackage{amsmath,amssymb}
\usepackage{graphicx}
\usepackage{xcolor}
\usepackage{tcolorbox}
\tcbuselibrary{breakable}
\usepackage{enumitem}

\newtcolorbox{infobox}[1]{colback=blue!5,colframe=blue!75!black,fonttitle=\bfseries,title=#1,breakable}
\newtcolorbox{examplebox}[1]{colback=green!5,colframe=green!50!black,fonttitle=\bfseries,title=#1,breakable}
\newtcolorbox{exercicebox}[1]{colback=orange!5,colframe=orange!75!black,fonttitle=\bfseries,title=#1,breakable}

\title{\textbf{Chapitre 1 : Représentation Numérique}\\
\large Exercices d'Application (à faire pendant le cours)}
\author{Dr. El Hadji Bassirou TOURÉ\\
Département de Mathématiques et Informatique\\
Université Cheikh Anta Diop de Dakar}
\date{}

\begin{document}
\maketitle

\begin{infobox}{Objectif}
Ces exercices sont à réaliser \textbf{pendant le cours}, immédiatement après chaque notion. Ils permettent de vérifier la compréhension avant de passer aux TDs.
\end{infobox}

%==============================================================================
\section{Exercice A1 : Interprétation Signed/Unsigned}
%==============================================================================

\begin{exercicebox}{Énoncé}
Soit le motif binaire $10110100_2$ sur 8 bits. Calculer sa valeur décimale :
\begin{enumerate}[label=\alph*)]
    \item Interprété comme \textbf{unsigned} (non signé)
    \item Interprété comme \textbf{signed} (signé, complément à deux)
\end{enumerate}
\end{exercicebox}

\vspace{0.5em}

\begin{center}
\includegraphics[width=0.95\textwidth]{fig_exo1_signed_unsigned.pdf}
\end{center}

\begin{examplebox}{Méthode}
\begin{enumerate}
    \item \textbf{Unsigned (B2U)} : Additionner les puissances de 2 où le bit vaut 1
    \[B2U_8 = \sum_{i=0}^{7} b_i \cdot 2^i\]
    \item \textbf{Signed (B2T)} : Si MSB = 1, le nombre est négatif
    \[B2T_8 = -b_7 \cdot 2^7 + \sum_{i=0}^{6} b_i \cdot 2^i = -128 + \text{(valeur des 7 bits restants)}\]
\end{enumerate}
\end{examplebox}

\vspace{1em}

\begin{exercicebox}{À vous !}
Appliquer la même méthode au motif $11100011_2$ sur 8 bits.

\vspace{0.5em}
\textbf{Réponses attendues :}
\begin{itemize}
    \item Unsigned : \rule{2cm}{0.4pt}
    \item Signed : \rule{2cm}{0.4pt}
\end{itemize}
\end{exercicebox}

\newpage
%==============================================================================
\section{Exercice A2 : Négation en Complément à Deux}
%==============================================================================

\begin{exercicebox}{Énoncé}
Calculer la négation de $x = 37$ sur 8 bits en utilisant la formule :
\[\boxed{-x = \sim x + 1}\]
(inverser tous les bits, puis ajouter 1)
\end{exercicebox}

\vspace{0.5em}

\begin{center}
\includegraphics[width=0.95\textwidth]{fig_exo2_negation.pdf}
\end{center}

\begin{examplebox}{Méthode}
\begin{enumerate}
    \item Écrire $x$ en binaire sur 8 bits
    \item Inverser tous les bits (0 → 1, 1 → 0) pour obtenir $\sim x$
    \item Ajouter 1 au résultat
    \item \textbf{Vérification} : $x + (-x)$ doit donner 0
\end{enumerate}
\end{examplebox}

\vspace{1em}

\begin{exercicebox}{À vous !}
Calculer la négation de $x = 100$ sur 8 bits.

\vspace{0.5em}
\begin{tabular}{ll}
$x = 100$ en binaire : & \rule{4cm}{0.4pt} \\[0.5em]
$\sim x$ : & \rule{4cm}{0.4pt} \\[0.5em]
$\sim x + 1 = -x$ : & \rule{4cm}{0.4pt} \\[0.5em]
Vérification $100 + (-100)$ : & \rule{4cm}{0.4pt}
\end{tabular}
\end{exercicebox}

\newpage
%==============================================================================
\section{Exercice A3 : Encodage IEEE 754}
%==============================================================================

\begin{exercicebox}{Énoncé}
Encoder le nombre $-10.75$ en format IEEE 754 float (32 bits).

Donner le résultat final en hexadécimal.
\end{exercicebox}

\vspace{0.5em}

\begin{center}
\includegraphics[width=0.95\textwidth]{fig_exo3_ieee754.pdf}
\end{center}

\begin{examplebox}{Méthode (6 étapes)}
\begin{enumerate}
    \item \textbf{Signe} : négatif → $s = 1$, positif → $s = 0$
    \item \textbf{Conversion binaire} : convertir $|x|$ en binaire (partie entière + fractionnaire)
    \item \textbf{Normalisation} : mettre sous forme $1.xxxxx \times 2^E$ (décaler la virgule)
    \item \textbf{Exposant biaisé} : calculer $E + 127$, puis convertir en binaire (8 bits)
    \item \textbf{Fraction} : prendre les bits après le "1." (23 bits, le 1 est implicite)
    \item \textbf{Assembler} : concaténer signe | exposant | fraction
\end{enumerate}
\end{examplebox}

\vspace{1em}

\begin{exercicebox}{À vous !}
Encoder $+6.5$ en IEEE 754 float (32 bits).

\vspace{0.5em}
\begin{tabular}{ll}
① Signe : & $s = $ \rule{1cm}{0.4pt} \\[0.5em]
② Binaire de $6.5$ : & \rule{4cm}{0.4pt} \\[0.5em]
③ Normalisation : & $1.$ \rule{2cm}{0.4pt} $\times 2^{E}$ où $E = $ \rule{1cm}{0.4pt} \\[0.5em]
④ Exposant biaisé : & $E + 127 = $ \rule{1cm}{0.4pt} $ = $ \rule{3cm}{0.4pt} (binaire) \\[0.5em]
⑤ Fraction (23 bits) : & \rule{5cm}{0.4pt} \\[0.5em]
⑥ Résultat hex : & \texttt{0x}\rule{2cm}{0.4pt}
\end{tabular}
\end{exercicebox}

\end{document}